% page 544 of the facsimile (image p-543 of the scan)
% block 160.0 x 242.0 mm ; left margin 8.0 mm
\pgc{-12.700mm}{0.000mm}{
\l{\cel{3}536}
\l{}
\l{\sou{sketar}. (\sou{netrans.}) Glite kurar sur glacio-surfaco, +werante}
\l{\cel{3}sorto de sandalo di qua la subajo esas provizita ye fero-}
\l{\cel{3}lamo vertikala. - E.}
\l{}
\l{\sou{skio}. Longa sketilo ek ligno, uzata por glitar sur nivo-sur-}
\l{\cel{3}faco. - DEFIS.}
\l{}
\l{\sou{skiro}. (\sou{patol.)}\cel{1}Tumoro sen-dolora, qua naskas precipue en la}
\l{\cel{3}korpo-parti glandatra, ulakaze kancera. -}
\l{}
\l{\sou{skisar}. (\sou{trans.)}\cel{2}I. Trasar, per krayono, per kreto, e c. la}
\l{\cel{3}linei precipua di desegnoto, di piktoto. - II. Facar la mo-}
\l{\cel{3}delo vaxa, argilo, e c. di skulturo. - III. Trasar la \textquotedbl{}mapo\textquotedbl{}}
\l{\cel{3}komencala di verko literaturala o ciencala. - IV. Deskriptar}
\l{\cel{3}per (lua) traiti precipua.- DFIRS.}
\l{}
\l{\sou{skism}o. I. En religio establisita : formaco di eklezio qua su}
\l{\cel{3}separas de la eklezio til nun agnoskita, sen disidenteso}
\l{\cel{3}kompleta pri la punti precipua di la dogmato o di la kulto.-}
\l{\cel{3}II. Su-separo en la domeni \textquotedbl{}politiko\textquotedbl{}, \textquotedbl{}arti\textquotedbl{}, e c. - DEFIRS}
\l{}
\l{\sou{skisto.}\cel{2}(mineral.) Roko stratoza, exfoliebla ye lameli dina,}
\l{\cel{3}e di.qua la elementi precipua esas silico, argilo, alumino,}
\l{\cel{3}o metal-oxidi diversa. - EFIS.}
\l{}
\l{\sou{sklavo}. I. La homo qua, en socio, ne esas en la strato de la}
\l{\cel{3}liberi. - II. (a) Ta qua su submisas ad ulu, obedias ulu,}
\l{\cel{3}sen rezisto, pro ambiciar ulo, por obtenar ulo ; (b) Ta qua}
\l{\cel{3}su submisas a la dominaco da muliero, pro amorar elu.-}
\l{\cel{3}DEFIS.}
\l{}
\l{\sou{sklerotik}a. (\sou{anat.}) Membrano blanka ye qua konsistas parto tre}
\l{\cel{3}granda di la exterajo di la okul-globo. - DEFIS.}
\l{}
\l{\sou{skol}o. I. (a) Establisuro en qua onu docas a la pueri la rudi-}
\l{\cel{3}menti : lekto, skribo, kalkulo, gramatiko elementala, e c.;}
\l{\cel{3}(b) Establisuro ube onu docas specalaji. - II. Ensemblo ek}
\l{\cel{3}la dicipuli qui studias ta o ca doktrino (filozofiala, me-}
\l{\cel{3}dicinala, artala, literaturala) - ed anke ta doktrino ipsa.}
\l{\cel{3}- III. To quo esas apta igar posedar la experienco di ulo}
\l{\cel{3}(la \textquotedbl{}skolo di vivo\textquotedbl{}, e c.). - DEFIRS.}
\l{}
\l{skolastiko.\cel{2}La filozofio quan onu docis en la skoli di la}
\l{\cel{3}mezepoko. - DEFIRS.}
\l{}
\l{\sou{skol}io. I. (\sou{gram.}) Noturo da komentero antiqua od olima. - II.}
\l{\cel{3}(geom.) Expliko pri teoremo, por igar lua apliko, plu kom-}
\l{\cel{3}pleta o limitoza. - DEFIS.}
\l{}
\l{\sou{skoliasto}. Komentero (antiqua od olima) di la texti Greka o}
\l{\cel{3}Latina. - DEFIRS.}
\l{}
\l{\sou{skolopendro}. (\sou{zool.)}\cel{2}Genero de miriapodo, insekto. - L.\cel{1}\sou{sco}-}
\l{\cel{3}\sou{lopendra}.- DEFIS.}
}
